% Part: sets-functions-relations
% Chapter: functions
% Section: partial-functions

\documentclass[../../../include/open-logic-section]{subfiles}


\begin{document}

\olfileid{sfr}{fun}{par}

\olsection{Partiële functies: niet altijd een uitkomst}

\begin{explain}
Soms maken we de definitie van een functie wat ruimer. Dan hoeft niet
elke mogelijke invoer een uitkomst te hebben. Zo'n functie heet een
\emph{partiële functie}.
\end{explain}

\begin{defn}
Een \emph{partiële functie} $f \colon A \pto B$ koppelt aan elk
!!{element} van~$A$ hoogstens één !!{element} van~$B$. Een invoer kan
dus zonder uitkomst blijven, maar krijgt nooit meer dan één uitkomst. Koppelt $f$ een
element van~$B$ aan $x \in A$, dan is $f(x)$ \emph{gedefinieerd}.
Anders noemen we dit \emph{ongedefinieerd}. Als $f(x)$ gedefinieerd
is, schrijven we $f(x) \fdefined$; anders schrijven we $f(x)
\fundefined$. Het \emph{domein} van een partiële functie~$f$ bestaat
uit de invoeren in~$A$ waarvoor wel een uitkomst bestaat. Dus
$\dom{f} = \Setabs{x \in A}{f(x) \fdefined}$.
\end{defn}

\begin{ex}
Elke functie $f\colon A \to B$ is ook een partiële functie. Een
partiële functie die voor elke invoer uit~$A$ een uitkomst heeft---dus
precies wat we tot nu toe gewoon een functie noemden---heet ook een
\emph{totale} functie.
\end{ex}

\begin{ex}
De partiële functie $f \colon \Real \pto \Real$ met $f(x) = 1/x$ heeft
geen uitkomst bij $x = 0$ en is voor elke andere invoer wel
gedefinieerd.
\end{ex}

\begin{prob}
Gegeven is $f\colon A \pto B$. Definieer de partiële functie $g\colon
B \pto A$ als volgt. Neem een $y \in B$. Is er precies één $x \in A$
waarvoor $f(x) = y$, zet dan $g(y) = x$. Schrijf in alle andere
gevallen $g(y) \fundefined$. Toon nu aan: als $f$ injectief is, dan
$g(f(x)) = x$ voor alle $x \in \dom{f}$, en $f(g(y)) = y$ voor alle
$y \in \ran{f}$.
\end{prob}

\begin{defn}[Grafiek van een partiële functie]
Neem een partiële functie $f\colon A \pto B$. De \emph{grafiek} van~$f$
houdt bij welke invoer welke uitkomst krijgt. Het is de relatie
$R_f \subseteq A \times B$ die wordt gegeven door
\[
R_f = \Setabs{\tuple{x,y}}{f(x) = y}.
\]
\end{defn}

\begin{prop}
Neem een relatie $R \subseteq A \times B$ met deze regel: als $Rxy$ en
$Rxy'$ allebei gelden, dan $y = y'$. Eén invoer kan dus niet twee
verschillende uitkomsten hebben. Dan is $R$ de grafiek van de partiële
functie $f\colon A \pto B$ die zo werkt: is er een $y$ waarvoor $Rxy$,
dan zetten we $f(x) = y$; anders schrijven we $f(x) \fundefined$. Is
$R$ ook \emph{serieel}, dan heeft elke invoer een uitkomst: voor elke
$x \in A$ is er een $y \in B$ waarvoor $Rxy$. In dat geval is $f$
totaal.
\end{prop}

\begin{proof}
Stel dat er een $y$ is waarvoor $Rxy$. Er kan geen andere $y' \neq y$
zijn waarvoor $Rxy'$, want dan voldoet $R$ niet aan de regel hierboven.
Als er een $y$ is waarvoor $Rxy$, is die $y$ dus uniek en geeft $f$ bij
die invoer een vaste uitkomst. Verder is $R_f = R$. Ook is $f$ totaal
als~$R$ serieel is.
\end{proof}
\end{document}
